\documentclass[11pt,twocolumn]{article}

\usepackage[margin=1in]{geometry}
\usepackage{amsmath,amssymb,amsthm}
\usepackage{graphicx}
\usepackage{hyperref}
\usepackage{listings}
\usepackage{xcolor}
\usepackage{booktabs}
\usepackage{enumitem}

\hypersetup{colorlinks=true,linkcolor=blue!60!black,citecolor=green!40!black,urlcolor=blue!60!black}

\lstset{
  language=Rust,
  basicstyle=\ttfamily\scriptsize,
  keywordstyle=\color{blue!70!black}\bfseries,
  commentstyle=\color{green!50!black}\itshape,
  stringstyle=\color{red!60!black},
  breaklines=true,
  frame=single,
  backgroundcolor=\color{gray!5}
}

\newtheorem{definition}{Definition}
\newtheorem{theorem}{Theorem}

\title{Post-Quantum Decentralized Identity: DID Registry, MFA, and Trust Credentials on Holochain}
\author{Tristan Stoltz \\ Luminous Dynamics \\ \texttt{tristan.stoltz@evolvingresonantcocreationism.com}}
\date{March 2026}

\begin{document}

\twocolumn[
\begin{@twocolumnfalse}
\maketitle

\begin{abstract}
We present a decentralized identity system implemented on Holochain with 13 zomes
spanning DID management, multi-factor authentication, trust credentials, verifiable
credentials, social recovery, and a web-of-trust. The system is distinguished by
three innovations: (1) a 5-level MFA assurance scheme that maps directly to the
identity dimension of a consciousness-gated governance profile; (2) post-quantum
key support via ML-KEM-768 alongside classical Ed25519, providing quantum resistance
without sacrificing current interoperability; and (3) a web-of-trust zome that computes
transitive trust for Sybil detection and community trust scoring. The identity cluster
serves as the foundation for a 16-cluster civic operating system, with 100+ sweettest
integration tests and 23+ unit tests validating the complete identity lifecycle from
DID creation through credential issuance, social recovery, and revocation.
\end{abstract}
\end{@twocolumnfalse}
]

\section{Introduction}

Identity is the foundation of governance. Every governance action---voting,
proposing, delegating, judging---requires knowing \emph{who} is acting and
\emph{how much} to trust them. Centralized identity systems (government IDs,
corporate SSO) concentrate this power in institutions that can revoke identity
at will. Self-sovereign identity (SSI) aims to return control to the
individual~\cite{tobin2016}.

Existing SSI implementations face several challenges:

\begin{enumerate}[nosep]
  \item \textbf{Blockchain dependency}: Most DID methods (did:ethr, did:ion)
    depend on blockchain infrastructure with its associated costs and latency.
  \item \textbf{Quantum vulnerability}: Classical cryptographic keys (RSA,
    Ed25519) are vulnerable to quantum computers~\cite{shor1994}. Most SSI
    systems have no quantum migration path.
  \item \textbf{Trust bootstrapping}: SSI systems provide identity but not
    trust. A verifiable credential proves an attribute but not whether the
    issuer should be trusted.
  \item \textbf{Recovery}: Key loss in SSI typically means permanent identity
    loss. Social recovery mechanisms are rare.
\end{enumerate}

We present the Mycelix Identity cluster: a 13-zome Holochain implementation
that addresses all four challenges.

\section{Architecture}

\subsection{The did:mycelix Method}

Each agent registers a DID using the method-specific scheme:

\begin{lstlisting}[language={}]
did:mycelix:<base58-encoded-pubkey>
\end{lstlisting}

The DID document, stored on the agent's source chain and replicated to the
DHT, contains:

\begin{lstlisting}
{
  "id": "did:mycelix:5GrwvaEF...",
  "authentication": [{
    "type": "Ed25519VerificationKey2020",
    "publicKeyMultibase": "z6Mk..."
  }],
  "keyAgreement": [{
    "type": "ML-KEM-768-2024",
    "publicKeyMultibase": "z7Mk..."
  }],
  "service": [{
    "type": "MycelixCluster",
    "serviceEndpoint": ["commons", "governance"]
  }]
}
\end{lstlisting}

\subsection{Zome Overview}

The 13 zomes form a layered architecture:

\textbf{Foundation layer:}
\begin{enumerate}[nosep]
  \item \texttt{did\_registry}: DID CRUD, key rotation, document resolution
  \item \texttt{mfa}: Multi-factor authentication with 5 assurance levels
  \item \texttt{recovery}: Social $m$-of-$n$ key recovery
  \item \texttt{revocation}: Key and credential revocation with DHT propagation
\end{enumerate}

\textbf{Credential layer:}
\begin{enumerate}[nosep,start=5]
  \item \texttt{trust\_credential}: Peer-to-peer trust attestations
  \item \texttt{verifiable\_credential}: W3C VC issuance and verification
  \item \texttt{credential\_schema}: Schema registry for credential types
  \item \texttt{education}: Educational credential management
\end{enumerate}

\textbf{Network layer:}
\begin{enumerate}[nosep,start=9]
  \item \texttt{name-registry}: Human-readable name resolution
  \item \texttt{web-of-trust}: Transitive trust computation
\end{enumerate}

\textbf{Integration:}
\begin{enumerate}[nosep,start=11]
  \item \texttt{bridge}: Cross-cluster identity dispatch
\end{enumerate}

\section{Multi-Factor Authentication}

\subsection{Five Assurance Levels}

The MFA zome defines five assurance levels that map directly to the identity
dimension ($I$) of the consciousness profile:

\begin{table}[h]
\centering
\begin{tabular}{llc}
\toprule
\textbf{Level} & \textbf{Requirements} & $I$ \\
\midrule
Anonymous & Key only & 0.00 \\
Basic & Email/phone & 0.25 \\
Verified & 2+ factors & 0.50 \\
HighlyAssured & HW key + biometric & 0.75 \\
Critical & In-person proofing & 1.00 \\
\bottomrule
\end{tabular}
\caption{MFA assurance levels and identity scores.}
\end{table}

Each assurance level is \emph{additive}: an agent can upgrade their level
by adding verification factors without re-verifying existing ones. The MFA
zome maintains an attestation chain linking each factor to the DID.

\subsection{Sybil Cost Curve}

The assurance levels create an exponential Sybil cost curve. An attacker
creating $n$ Sybil identities at level $l$ faces cost:

\begin{equation}
  \text{Cost}(n, l) = n \cdot c(l)
\end{equation}

where $c(\text{Anonymous}) \approx 0$ (just a keypair), $c(\text{Basic}) \approx$
cost of a phone number, and $c(\text{Critical}) \approx$ cost of in-person
social engineering of trusted community members.

At the Critical level, creating 10 Sybil identities requires convincing 10
separate groups of trusted community members to perform in-person identity
proofing---a cost that scales linearly with $n$ and is bounded by the
attacker's social capacity.

\section{Post-Quantum Cryptography}

\subsection{Hybrid Key Architecture}

Each DID document can contain both classical and post-quantum keys:

\begin{description}[nosep]
  \item[Authentication] Ed25519 (classical) --- used for source chain signing,
    DHT entry validation, and all current Holochain operations.
  \item[Key Agreement] ML-KEM-768 (post-quantum) --- used for encrypted
    communication channels, future-proofing against quantum adversaries.
\end{description}

The hybrid approach ensures:
\begin{itemize}[nosep]
  \item Current compatibility with Holochain's Ed25519 infrastructure
  \item Quantum resistance for key agreement (encrypted mail, private channels)
  \item Forward secrecy: messages encrypted today remain secure against future
    quantum computers
\end{itemize}

\subsection{ML-KEM-768 Integration}

ML-KEM-768 (Module Lattice Key Encapsulation Mechanism)~\cite{nist2024pqc}
was standardized by NIST in 2024 as FIPS 203. Key properties:

\begin{itemize}[nosep]
  \item Security level: AES-192 equivalent (NIST Category 3)
  \item Public key size: 1,184 bytes
  \item Ciphertext size: 1,088 bytes
  \item Encapsulation/decapsulation: $<1$ ms
\end{itemize}

The larger key sizes (compared to 32-byte Ed25519 keys) are stored in the
DID document on the agent's source chain. DHT replication propagates the
full document to the storage neighborhood.

\subsection{Key Rotation Protocol}

When an agent upgrades from classical-only to hybrid (classical + PQ) keys:

\begin{enumerate}[nosep]
  \item Generate ML-KEM-768 key pair locally
  \item Create a DID document update entry signed by the existing Ed25519 key
  \item The update includes the new ML-KEM-768 public key in \texttt{keyAgreement}
  \item DHT validators verify the Ed25519 signature on the update
  \item The new document replaces the old in the DHT neighborhood
\end{enumerate}

This protocol ensures that quantum key addition is authenticated by the
classical key, maintaining the chain of trust.

\section{Trust Credentials}

\subsection{Peer-to-Peer Attestation}

The \texttt{trust\_credential} zome enables agents to attest to each other's
trustworthiness. A trust attestation contains:

\begin{lstlisting}
pub struct TrustAttestation {
    pub subject: AgentPubKey,
    pub attestor: AgentPubKey,
    pub trust_level: f64,  // [0, 1]
    pub domain: Option<String>,
    pub evidence: Option<String>,
    pub created_at: Timestamp,
}
\end{lstlisting}

Trust attestations feed the community dimension ($C$) of the consciousness
profile. The aggregation weights attestations by the attestor's own
consciousness tier:

\begin{equation}
  C = \frac{\sum_{j} \tau_j \cdot w(\text{tier}_j)}{\sum_{j} w(\text{tier}_j)}
\end{equation}

This prevents Sybil identities (all at Observer tier, weight 0) from
bootstrapping each other's community scores.

\subsection{Verifiable Credentials}

The \texttt{verifiable\_credential} zome implements W3C Verifiable
Credentials~\cite{w3cvc} with Holochain-native verification:

\begin{enumerate}[nosep]
  \item Credentials are stored on the holder's source chain (holder-controlled)
  \item Verification uses the DHT to resolve the issuer's DID and public key
  \item Selective disclosure is supported via derived credentials
  \item Credential schemas are registered in the \texttt{credential\_schema} zome
\end{enumerate}

\section{Social Recovery}

\subsection{$m$-of-$n$ Recovery Protocol}

The \texttt{recovery} zome implements social key recovery:

\begin{enumerate}[nosep]
  \item An agent designates $n$ \emph{recovery guardians} (trusted peers)
  \item The agent's recovery secret is split into $n$ shares using Shamir's
    Secret Sharing
  \item Each guardian stores one share on their source chain
  \item To recover: $m$ guardians provide their shares to reconstruct the
    recovery secret
  \item The recovery secret authorizes a key rotation on the DID document
\end{enumerate}

\begin{definition}[Recovery threshold]
The recovery threshold $m$ is chosen such that $m > n/2$ (majority required)
and $m < n$ (not all guardians needed). A typical configuration is $m = 3$
of $n = 5$.
\end{definition}

\subsection{Guardian Selection}

Guardian selection is guided by the web-of-trust: agents with higher trust
scores are recommended as guardians. The system warns if:

\begin{itemize}[nosep]
  \item Selected guardians are too closely connected (correlated failure risk)
  \item Selected guardians have low consciousness tiers (low reliability)
  \item The $m/n$ ratio is too low (easy to compromise) or too high (hard to
    recover)
\end{itemize}

\section{Web of Trust}

\subsection{Transitive Trust Computation}

The \texttt{web-of-trust} zome computes transitive trust across the network.
Given direct trust edges $t(A, B) \in [0, 1]$, the transitive trust from
$A$ to $C$ through $B$ is:

\begin{equation}
  t(A \to B \to C) = t(A, B) \cdot t(B, C) \cdot d
\end{equation}

where $d \in (0, 1)$ is the per-hop decay factor (default 0.8). Multi-path
trust uses the maximum across all paths up to a configurable depth limit.

\subsection{Sybil Detection}

The web-of-trust enables structural Sybil detection. A cluster of Sybil
identities creates a dense subgraph with few connections to the legitimate
network. The web-of-trust zome identifies such clusters by computing:

\begin{equation}
  \text{isolation}(G) = 1 - \frac{|\text{edges}(G, \bar{G})|}{|G| \cdot |\bar{G}|}
\end{equation}

where $G$ is the suspected cluster and $\bar{G}$ is the rest of the network.
High isolation scores flag potential Sybil clusters for community review.

\section{Name Registry}

The \texttt{name-registry} zome provides human-readable name resolution:

\begin{lstlisting}[language={}]
alice.mycelix -> did:mycelix:5Grwva...
cooperative-housing.mycelix -> did:mycelix:7Xbf...
\end{lstlisting}

Names are first-come-first-served with consciousness tier requirements
(Participant+ to register, Steward+ to claim premium names). Name disputes
use the Justice domain in the Civic cluster.

\section{Revocation}

The \texttt{revocation} zome handles two types of revocation:

\begin{description}[nosep]
  \item[Key revocation] Marks a key as compromised. The DID document is updated
    with a new key, and the revocation entry propagates through the DHT.
  \item[Credential revocation] Marks a specific credential as invalid. Verifiers
    check the revocation status during credential verification.
\end{description}

Revocation is immediate on the agent's source chain and eventually consistent
across the DHT (typically seconds to minutes depending on network size).

\section{Integration with Consciousness Gating}

The Identity cluster feeds two dimensions of the consciousness profile:

\begin{enumerate}[nosep]
  \item \textbf{Identity dimension} ($I$): Directly from the MFA assurance level
  \item \textbf{Community dimension} ($C$): Aggregated from trust credential
    attestations, weighted by attestor tier
\end{enumerate}

The bridge zome provides \texttt{get\_consciousness\_credential}, which queries
the MFA zome for the current assurance level, the trust credential zome for
community attestations, and the reputation bridge for behavioral history, then
constructs a \texttt{SovereignProfile} (or legacy \texttt{ConsciousnessProfile} for backward compatibility) and issues a 24-hour
\texttt{SovereignCredential}.

\section{Evaluation}

\subsection{Test Coverage}

\begin{itemize}[nosep]
  \item 23+ unit tests covering DID lifecycle, MFA level transitions, trust
    aggregation, and recovery protocol
  \item 100+ sweettest integration tests with real Holochain conductor,
    covering multi-agent scenarios, DHT propagation, cross-zome calls
  \item Governance audit: 7 previously missing \texttt{\#[ignore]} annotations
    fixed for tests requiring conductor
\end{itemize}

\subsection{Key Size Trade-offs}

ML-KEM-768 public keys (1,184 bytes) are 37$\times$ larger than Ed25519 keys
(32 bytes). For source chain storage, this is negligible (source chains routinely
store multi-KB entries). For DHT bandwidth, the impact is proportional to
replication factor---with 3$\times$ replication, each DID document update
transmits $\sim$3.5 KB of key material.

\section{Related Work}

\textbf{did:web}~\cite{didweb}: DNS-based DID resolution. Simple but centralized
(DNS dependency). No quantum resistance.

\textbf{did:ion} (Sidetree/Bitcoin): Bitcoin-anchored DIDs. Quantum-vulnerable
and depends on Bitcoin's continued operation.

\textbf{did:key}: Ephemeral DIDs from public keys. No recovery, no MFA, no
trust network.

\textbf{Spruce/SpruceID}: SSI toolkit with multiple DID methods. Provides
credential infrastructure but not the integrated web-of-trust or
consciousness gating.

\section{Conclusion}

The Mycelix Identity cluster demonstrates that decentralized identity can be
both comprehensive (13 zomes covering the full identity lifecycle) and
future-proof (ML-KEM-768 post-quantum keys). The integration with
consciousness-gated governance---where identity verification strength directly
determines governance capability---creates a natural incentive for strong
identity practices without coercion.

The web-of-trust and social recovery mechanisms address the two most common
criticisms of self-sovereign identity: ``how do I know who to trust?'' and
``what if I lose my keys?'' By embedding these solutions in the same DHT
that hosts the identity records, Mycelix provides answers that are
decentralized, privacy-preserving, and operationally practical.

\begin{thebibliography}{99}
  \bibitem{tobin2016} Tobin, A. and Reed, D. (2016). The Inevitable Rise of
    Self-Sovereign Identity. Sovrin Foundation.
  \bibitem{shor1994} Shor, P. W. (1994). Algorithms for Quantum Computation.
    \textit{FOCS}.
  \bibitem{nist2024pqc} NIST (2024). Module-Lattice-Based Key-Encapsulation
    Mechanism Standard. FIPS 203.
  \bibitem{w3cdid} W3C (2022). Decentralized Identifiers (DIDs) v1.0.
  \bibitem{w3cvc} W3C (2022). Verifiable Credentials Data Model v1.1.
  \bibitem{didweb} did:web Method Specification. W3C CCG.
  \bibitem{brock2018} Brock, A. and Harris-Braun, E. (2018). Holochain:
    Scalable Agent-Centric Distributed Computing.
  \bibitem{shamir1979} Shamir, A. (1979). How to Share a Secret.
    \textit{Communications of the ACM}, 22(11).
\end{thebibliography}

\end{document}
